\scanpage{52}
\mspage{58}

d'où le diagramme cubique

\diagram{\begin{tikzcd}[ampersand replacement=\&,column sep=small,row sep=small]
X \arrow[dd] \& {} \& X' \arrow[ll] \arrow[dd] \& {} \\
{} \& X_{/y} \arrow[ul] \& {} \& X'_{/y'} \arrow[ul] \arrow[ll,"\unsure{k|\kappa}"'] \arrow[dd] \\
Y \& {} \& Y' \arrow[ll] \& {} \\
{} \& Y_{/y} \arrow[ul] \& {} \& Y'_{/y'} \arrow[ul] \arrow[ll,"\pi"] \arrow[from=2-2,to=4-2,crossing over]
\end{tikzcd}}

, où les \deleted{\unsure{côtés|cotés}} carrés latéraux (verticaux) sont cartésiens. On note que $\pi$ est lisse et à fibres \deleted{\unsure{car}}contractiles (\unsure{car|si} $p$ est lisse) donc $k$ aussi. \deleted{Notons \unsure{donc|alors}} Donc $k$ est une hot-équivalence \deleted{\unsure{stricte}} \deleted{\unsure{Cor.}} \deleted{\unsure{les}} \unsure{$=$} costricte \unsure{$g$|\}}, donc $\Gamma(X_{/y})\mathrel{\overset{\sim}{\unsure{\relbar\joinrel\relbar|\longrightarrow}}}\Gamma(X'_{/y'})$, on gagne encore.

\uline{NB} Sauf erreur, on a une réciproque du th., savoir : Si $f:X\longrightarrow Y$ est tel que le th. de chgt. de base pour $f_*(F)$ est vrai pour \uline{tout} dérivateur (il suffit sans doute de $\mathrm{HOTAB}$ et $\mathrm{HOT}$) et pour tout $Y'\unsure{\relbar\joinrel\relbar|\longrightarrow}Y$, alors $f$ est hom$^{\unsure{\mathrm t}}$ propre.

D'autre part, si $g:Y'\longrightarrow Y$ est tel que le th. de chang$^{\mathrm t}$ de base par $g$ est vrai pour tout $X\xrightarrow{f}Y$ et tout $F\in\unsure{\mathbb D|\underline D}(X)$ ($\underline D$ dérivateur \inserted{\unsure{pr. à dr.|gr. à dr.}} quelc.), alors $g$ est hom$^{\mathrm t}$ lisse.

Du côté homologique : les foncteurs \uline{lisses} $f$ sont caractérisés comme ceux pour lesquels \deleted{\unsure{R}} $f_!$ commute à tout changement de base. Les foncteurs propres $g$ sont exactement ceux pour lesquels le chang$^{\mathrm t}$ de base par $g$ commute à $f_!$, pour \deleted{\unsure{ex. :}} \inserted{tout} $f:X\longrightarrow Y$ et tout $\underline D$ \unsure{dér}
% archive sheet number: 51
